% proof_IV.tex -- LaTeX form of proof_IV.md.
% -----------------------------------------------------------------------------
% Standalone: compile with pdflatex. Nothing here is \input by main.tex; the
% theorem bodies are written to match main.tex's Section IV verbatim so blocks
% can be lifted across without re-deriving anything.
%
% Tags on each statement:
%   [PROVED]      complete argument, no unverified hypothesis
%   [ASSUMPTION]  needed, stated explicitly, measured violation rate given
%   [MEASURED]    a number from simulation, with the producing script named
%   [NOT CLAIMED] deliberately withheld, with the reason
% -----------------------------------------------------------------------------
\documentclass[11pt]{article}
\usepackage[margin=1in]{geometry}
\usepackage{amsmath,amssymb,amsthm}
\usepackage{booktabs}
\usepackage{hyperref}

\newtheorem{theorem}{Theorem}
\newtheorem{lemma}[theorem]{Lemma}
\newtheorem{proposition}[theorem]{Proposition}
\newtheorem{corollary}[theorem]{Corollary}
\theoremstyle{definition}
\newtheorem{definition}[theorem]{Definition}
\newtheorem{assumption}[theorem]{Assumption}
\theoremstyle{remark}
\newtheorem{remark}[theorem]{Remark}

\newcommand{\JEA}{J_{\mathrm{EA}}}
\newcommand{\JLQR}{J_{\mathrm{LQR}}}
\newcommand{\Kd}{K_{\mathrm{d}}}
\newcommand{\Ks}{K_{\mathrm{s}}}
\newcommand{\gf}{\gamma_{\mathrm{f}}}
\newcommand{\Rset}{\mathcal{R}}
\newcommand{\Sstar}{S^{\star}}
\newcommand{\tagged}[1]{\textsc{[#1]}}

\title{Section IV: proof repository\\[2pt]
\large Actuator co-design --- where the search stops and how far that is from optimal}
\author{}
\date{}

\begin{document}
\maketitle
\vspace{-2.5em}

\section{Setup}

Hold the link count $\ell$ and the sensor mask $\mathbf{s}$ fixed. Write
$\Rset := \{1,\dots,N_u\}$ for the actuator index set and, for a retained set
$S \subseteq \Rset$, let $K(S)$ denote $\Pi_\ell(\Kd)$ with rows outside $S$ and
columns outside $\mathbf{s}$ zeroed, so $K(\Rset) = \Ks([\ell,\mathbf{1},\mathbf{s}])$.

\begin{align}
  f(S) &:= -\,\JLQR(K(S))\,/\,\JLQR(\Kd)
    \quad\in\ [-\infty,\,0),
    \label{eq:f}\\[2pt]
  m_u &:= \bigl|\operatorname{supp}\bigl([\Pi_\ell(\Kd)]_{u,:}\bigr) \cap \mathbf{s}\bigr|,
    \label{eq:mu}\\[2pt]
  M(S) &:= \sum_{u \in S}\Bigl( w_a\,\mathbf{1}\{m_u > 0\} + w_c\,m_u \Bigr),
    \label{eq:M}\\[2pt]
  U(S) &:= f(S) - M(S).
    \label{eq:U}
\end{align}

Larger $f$ is better; $f = -\infty$ exactly on the destabilizing sets. With
$\ell$ and $\mathbf{s}$ fixed, the sensor and constant terms of $\JEA$ are common
to all $S$, so minimizing $\JEA$ over the actuator mask is equivalent to
maximizing $U$.

\begin{lemma}[$M$ is exactly modular]\label{lem:modular} \tagged{proved}
$M$ is a sum of per-element weights fixed by $(\ell,\mathbf{s})$, hence modular.
\end{lemma}

\begin{proof}
Distinct rows of $K(S)$ have disjoint supports, so $N_c(K(S)) = \sum_{u\in S} m_u$
exactly. An actuator whose row is emptied by the truncation $\Pi_\ell$ carries no
controller entries and is not charged $w_a$ by $\JEA$, which is what the indicator
in \eqref{eq:M} records. Each summand therefore depends on $u$ alone.
\end{proof}

\begin{remark}[why the indicator is not optional]
Writing $w_a|S|$ in place of $w_a\,\#\{u \in S : m_u > 0\}$ overcharges every
masked-on actuator whose row the truncation emptied. The two coincide whenever
$\ell$ is large enough that every masked row survives, which is why the EA never
exposed the difference ($\ell$ at $90$--$99\%$ of its maximum) and the greedy
baseline immediately did ($\ell$ at $5$--$12\%$): at one greedy fixed point $16$
actuators were masked on but only $4$ had entries, inflating $M(\Sstar)$ from
$2.0$ to $6.8$ and the reported gap from $2.02$ to $6.82$. All of the curvature of
the problem sits in $f$; by Lemma \ref{lem:modular} the structural terms
contribute none.
\end{remark}

\section{Assumptions}

\begin{assumption}[approximate submodularity]\label{as:submod} \tagged{assumption}
Let $\mathcal{B}$ be the range of architecture sizes the search occupies. For all
$S$ with $|S| \in \mathcal{B}$ and $A+BK(S)$ Schur stable, and all $T$,
\begin{equation}\label{eq:submod}
  \sum_{u \in T}\bigl[f(S \cup u) - f(S)\bigr]
  \;\geq\; \gf\,\bigl[f(S \cup T) - f(S)\bigr],
  \qquad \gf > 0 .
\end{equation}
\end{assumption}

\begin{assumption}[monotonicity]\label{as:mono} \tagged{assumption}
On the same sets, $S \subseteq T \Rightarrow f(S) \leq f(T)$.
\end{assumption}

Neither is automatic: $K(S)$ is a truncation of $\Kd$, not the optimal gain for
the support $S$.

\paragraph{\tagged{measured} The band in Assumption \ref{as:submod} is not cosmetic.}
Sampling the Das--Kempe quotient \eqref{eq:submod} stratified by retention
fraction, $250$ draws per cell, $|T| \geq 2$ (with $|T| = 1$ the quotient is
identically $1$ and carries no information):

\begin{center}
\begin{tabular}{lcccc}
\toprule
$|S|/N_u$ & $0.05$--$0.15$ & $0.15$--$0.30$ & $0.30$--$0.45$ & $0.45$--$0.80$ \\
\midrule
$5\times5$ s1  & $1.000$ & $0.980$  & $0.981$  & $0.911$ \\
$5\times5$ s10 & $1.000$ & $0.999$  & $0.871$  & $0.879$ \\
$7\times7$ s1  & $0.976$ & $-0.119$ & $-2.012$ & $0.717$ \\
$7\times7$ s10 & $1.000$ & $0.986$  & $0.652$  & $0.110$ \\
\bottomrule
\end{tabular}
\end{center}

Entries are $\min$ over samples. The searches return $|\Sstar|$ of order $10\%$ of
$N_u$, i.e.\ the leftmost column, where $f$ is submodular. The negative values sit
in the middle bands, where the aggregate gain $f(S \cup T) - f(S)$ is near zero and
the quotient is dominated by its denominator; $18$--$26\%$ of draws there are
discarded for that reason. Estimating $\gf$ on the band the search occupies gives
$\gf = 1.0000$ on all three plants of the paper.

\paragraph{\tagged{measured} Assumption \ref{as:mono}.}
Over $600$ sampled $(S,u)$ pairs with $S$ feasible: $0$ violations on the
$5\times5$ grid, $2$ on the $7\times7$, of relative size $\leq 0.8\%$, and no
sampled addition destabilized. Script: \texttt{scratchpad/probe\_A2.m}.

\section{Where the search stops}

\begin{definition}[$1$-flip local optimum]\label{def:loc}
$\theta^\star$ is a \emph{$1$-flip local optimum} if no single actuator or sensor
mask flip strictly decreases $\JEA$.
\end{definition}

\begin{proposition}[elite-child hitting probability]\label{prop:hit} \tagged{proved}
Let $\mathcal{H}_{t-1}$ be the $\sigma$-algebra generated by the population history
through generation $t-1$. For every single-bit neighbour $\theta'$ of the elite,
\begin{equation}\label{eq:qmin}
  \mathbb{P}\bigl(\theta^{\mathrm{c}} = \theta' \mid \mathcal{H}_{t-1}\bigr)
  \;\geq\; q_{\min}
  \;:=\; \frac{(1-p_c)\,p_m(1-p_m)^{N_u+N_x-1}}{N_p\,(2d+1)}
  \;>\; 0
\end{equation}
uniformly in $t$, where $\theta^{\mathrm{c}}$ is some offspring of generation $t$.
\end{proposition}

\begin{proof}
With probability $1-p_c$ the algorithm skips crossover and copies the fitter of the
two selected parents. The elite has the strictly smallest cost, so whenever it is
selected it is the fitter one; selection assigns it weight at least $1/N_p$.
Mutation then flips exactly the required bit and no other with probability
$p_m(1-p_m)^{N_u+N_x-1}$, and leaves $\ell$ unchanged with probability $1/(2d+1)$.
\end{proof}

\begin{theorem}[almost-sure finite-time convergence]\label{thm:conv} \tagged{proved}
Suppose $\JEA(\theta^\star_0) < \infty$. Then almost surely there is a finite random
$T$ with $\JEA(\theta^\star_t) = \JEA(\theta^\star_T)$ for all $t \geq T$, and the
limiting gene is a $1$-flip local optimum.
\end{theorem}

\begin{proof}
$\Theta$ is finite, so $\JEA(\Theta)$ is a finite value set. Elitism carries the
$n_e$ best genes forward unchanged, so $\JEA(\theta^\star_t)$ is non-increasing; a
non-increasing sequence in a finite ordered set is eventually constant, at some
$J_\infty$. For $t > T$, Proposition \ref{prop:hit} gives each prescribed
single-bit neighbour probability at least $q_{\min} > 0$ per generation, and
$\sum_{t>T} q_{\min} = \infty$, so by the second Borel--Cantelli lemma every such
neighbour is produced infinitely often almost surely. If any had cost below
$J_\infty$, elitism would carry it forward, contradicting the definition of
$J_\infty$.
\end{proof}

\begin{proposition}[the returned architecture is irreducible]\label{prop:irr} \tagged{proved}
At the limiting gene, with $\Sstar$ its retained actuator set,
\begin{equation}\label{eq:irr}
\begin{aligned}
  u \in \Sstar &\;\Longrightarrow\; f(\Sstar) - f(\Sstar \setminus u) \;\geq\; w_a + w_c m_u, \\
  u \notin \Sstar &\;\Longrightarrow\; f(\Sstar \cup u) - f(\Sstar) \;\leq\; w_a + w_c m_u .
\end{aligned}
\end{equation}
\end{proposition}

\begin{proof}
Immediate from Definition \ref{def:loc} applied to $U = f - M$, using Lemma
\ref{lem:modular} for the marginal of $M$.
\end{proof}

Each retained actuator pays for itself and each discarded one does not. Both sides
of \eqref{eq:irr} are computable at a cost of $N_u - |\Sstar|$ Lyapunov solves,
which makes them usable as a stopping test.

\begin{remark}[what Theorem \ref{thm:conv} does not give]\label{rem:norate} \tagged{not claimed}
No rate. The chain on $\Theta$ is irreducible, so the classical result for elitist
EAs with full-support mutation would give almost-sure convergence to the global
optimum as $t \to \infty$; we do not state it, because it is asymptotic with no
rate and holds for \emph{any} elitist GA with full-support mutation.
\end{remark}

\section{Distance to optimality}

The search need not have reached a $1$-flip local optimum when the budget expires,
so the violation is carried explicitly. For $u \notin \Sstar$,
\begin{equation}\label{eq:slack}
  \xi_u := \bigl[\,f(\Sstar \cup u) - f(\Sstar) - w_a - w_c m_u\,\bigr]_+ ,
  \qquad
  \Xi(T) := \sum_{u \in T} \xi_u ,
\end{equation}
so that
\begin{equation}\label{eq:slackineq}
  f(\Sstar \cup u) - f(\Sstar) \;\leq\; w_a + w_c m_u + \xi_u
\end{equation}
holds with no hypothesis, and $\xi_u = 0$ exactly when $u$ cannot be added
profitably.

\begin{theorem}[optimality gap]\label{thm:gap} \tagged{proved under \ref{as:submod}, \ref{as:mono}}
For any $S$ with $A+BK(S)$ Schur stable,
\begin{equation}\label{eq:gap}
  U(S) - U(\Sstar)
  \;\leq\; M(\Sstar \setminus S)
  \;+\; \bigl(\gf^{-1} - 1\bigr)\,M(S \setminus \Sstar)
  \;+\; \gf^{-1}\,\Xi(S \setminus \Sstar).
\end{equation}
\end{theorem}

\begin{proof}
Assumption \ref{as:submod} with $T = S \setminus \Sstar$, followed by
\eqref{eq:slackineq}, gives
\begin{equation*}
  f(\Sstar \cup S) - f(\Sstar)
  \;\leq\; \gf^{-1}\!\!\!\sum_{u \in S \setminus \Sstar}\!\!\!
    \bigl[f(\Sstar \cup u) - f(\Sstar)\bigr]
  \;\leq\; \gf^{-1}\bigl[M(S \setminus \Sstar) + \Xi(S \setminus \Sstar)\bigr].
\end{equation*}
Assumption \ref{as:mono} gives $f(S) \leq f(\Sstar \cup S)$. Subtract $M(S)$, split
it modularly (Lemma \ref{lem:modular}) as
$M(S) = M(S \setminus \Sstar) + M(S \cap \Sstar)$, and use
$f(\Sstar) = U(\Sstar) + M(\Sstar \setminus S) + M(\Sstar \cap S)$.
\end{proof}

\begin{corollary}[computable bound]\label{cor:gap} \tagged{proved under \ref{as:submod}, \ref{as:mono}}
Since $\Sstar \setminus S \subseteq \Sstar$ and
$S \setminus \Sstar \subseteq \Rset \setminus \Sstar$,
\begin{equation}\label{eq:gapcomp}
  \min_{\mathbf{a}} \JEA \;\geq\; \JEA(\Sstar) - \mathrm{gap},
  \qquad
  \mathrm{gap} := M(\Sstar)
  + \bigl(\gf^{-1}\!-\!1\bigr) M(\Rset \setminus \Sstar)
  + \gf^{-1}\,\Xi(\Rset \setminus \Sstar),
\end{equation}
the minimum being over actuator masks at the current $(\ell,\mathbf{s})$. Every
term on the right is available from the run.
\end{corollary}

\paragraph{\tagged{measured} Both correction terms vanish in practice.}
With $\gf = 1$ on the operating band, \eqref{eq:gapcomp} collapses to
\begin{equation}\label{eq:gapsimple}
  \mathrm{gap} \;=\; M(\Sstar) \;+\; \Xi(\Rset \setminus \Sstar)
\end{equation}
--- the structural price of what was kept, plus what was left on the table by not
converging.

\begin{center}
\begin{tabular}{lccccc}
\toprule
plant, seed & $\gf$ & gap & $\JEA(\Sstar)$ & $\min_{\mathbf{a}}\JEA \geq$ & Thm.\ \ref{thm:gap} hypothesis \\
\midrule
$5\times5$, 1   & $1.0000$ & $2.100$ & $6.775$  & $4.675$ & $150/150$, end \checkmark \\
$5\times5$, 10  & $1.0000$ & $1.450$ & $7.811$  & $6.361$ & $63/150$,  end \checkmark \\
$5\times5$, 15  & $1.0000$ & $0.975$ & $6.547$  & $5.572$ & $24/150$,  end $\times$ \\
$7\times7$, 1   & $1.0000$ & $2.200$ & $7.646$  & $5.446$ & $144/150$, end \checkmark \\
$7\times7$, 10  & $1.0000$ & $2.248$ & $10.437$ & $8.189$ & $95/150$,  end $\times$ \\
$7\times7$, 15  & $1.0000$ & $7.199$ & $11.185$ & $3.986$ & $22/150$,  end $\times$ \\
IEEE 13-bus     & $1.0000$ & $0.450$ & $2.448$  & $1.998$ & $149/150$, end \checkmark \\
\bottomrule
\end{tabular}
\end{center}

Where $\Xi = 0$ the gap \emph{is} $M(\Sstar)$: $18$--$31\%$ of $\JEA$. The outlier,
$7\times7$ seed $15$ at $64\%$, is slack-dominated ($\Xi \approx 5.2$), and that
seed is $1$-flip locally optimal in only $22$ of $150$ generations. An unfinished
run receives a wider band, not a tighter unjustified one.

\begin{remark}[architecture size is set by the weights]\label{rem:size} \tagged{proved}
Every term of $\JEA$ is non-negative, so
\begin{equation}\label{eq:Nabound}
  w_a\,N_a(\Sstar) \;\leq\; \JEA(\Sstar)
  \qquad\Longrightarrow\qquad
  N_a(\Sstar) \;\leq\; \JEA(\Sstar)/w_a .
\end{equation}
The LQ term is normalized to $1$ at $\Kd$, so the entire LQ budget available to pay
for actuators is $O(1)$ while each costs $w_a$: the actuator count at the optimum
is set by the weights, not by the system size. Measured, $\JEA(\Sstar)/w_a$ is
$13$--$28$ against actual $|\Sstar|$ of $2$--$7$.
\end{remark}

\begin{remark}[scope]\label{rem:scope}
Three limits. \emph{(i)} \eqref{eq:gapcomp} bounds the optimum over actuator masks
at the elite's current $\ell$ and $\mathbf{s}$; relaxing those two coordinates can
buy more than the entire certified gap, so the band is not a bound on the global
optimum. Measured, the greedy architecture costs less than the bound on $2$ of $6$
runs --- no contradiction, since greedy varies $\ell$ and $\mathbf{s}$. The
restriction is forced: $f$ restricted to sensor sets is $-\infty$ on most subsets
and admits no submodularity ratio. \emph{(ii)} Theorem \ref{thm:gap} and
Proposition \ref{prop:irr} are statements about a \emph{point}, not an algorithm;
they hold for any method returning a $1$-flip local optimum. Only Proposition
\ref{prop:hit}, the proof of Theorem \ref{thm:conv}, and Remark \ref{rem:pm} are
specific to the EA. \emph{(iii)} $\gf$ and Assumption \ref{as:mono} are estimated,
not proved, so the results are rigorous conditional on constants we measure but
cannot certify.
\end{remark}

\begin{remark}[mutation-rate scaling]\label{rem:pm} \tagged{proved + measured}
The factor $(1-p_m)^{N_u+N_x-1}$ in \eqref{eq:qmin} decays exponentially in the
gene length at fixed $p_m$, but $p_m = c/(N_u+N_x)$ gives
$(1-p_m)^{N_u+N_x} \to e^{-c}$, so $q_{\min}$ becomes dimension-free. Measured,
$(1-1/n)^n = 0.3654$ at $n = 75$ and $0.3666$ at $n = 147$, against $0.0213$ and
$5.31\times10^{-4}$ for fixed $p_m = 0.05$. Ten-seed paired test of $p_m = 1/n$
against $p_m = 0.05$: $+1.56\%$ on the $5\times5$ grid ($5/10$ wins, $p = 0.183$),
$+11.36\%$ on the $7\times7$ ($9/10$ wins, $t = 3.37$, $p = 0.0083$). The predicted
signature is not ``scaled is better'' but ``the advantage grows with $n$'', and
that is what appears.
\end{remark}

\section{The greedy baseline}

Greedy round-robin local search over $\{\ell, \mathbf{a}, \mathbf{s}\}$, each phase
run to a fixed point, with the link phase restricted to the EA's own $\pm d$
neighbourhood.

\begin{proposition}[the descent path never leaves the stable set]\label{prop:G1} \tagged{proved}
Starting from $\Kd$, which stabilizes by construction, and accepting only strict
improvements, every iterate is feasible.
\end{proposition}

\begin{proof}
$\JEA = +\infty$ on infeasible points, so any infeasible candidate has cost
strictly greater than the current finite incumbent and is never accepted.
\end{proof}

\tagged{measured} Feasible on $6/6$ runs, while $4$--$10$ infeasible candidates are
probed and rejected ($215$--$251$ on the open-loop unstable plant). The EA has no
such property: mutation samples infeasible genes unconditionally, and it probes
$189$--$332$ of them per run. This contradicts the argument that removing a link
risks ``an infinite jump greedy cannot back out of'': on these plants one never has
to pass through the infeasible set at all.

\begin{proposition}[Theorem \ref{thm:gap} holds for greedy in its tight form]\label{prop:G3} \tagged{proved}
Greedy's stopping condition --- no single-bit flip improves $\JEA$ --- is strictly
stronger than the addition half used in \eqref{eq:slackineq}, so $\xi_u \equiv 0$
and \eqref{eq:gapcomp} reduces to $M(\Sstar) + (\gf^{-1}-1)M(\Rset\setminus\Sstar)$
by construction, irrespective of convergence.
\end{proposition}

\tagged{measured} Evaluated at the greedy fixed point the gap is $23$--$31\%$,
against $15$--$64\%$ at the EA's, whose spread is driven by $\Xi$.

\begin{proposition}[evaluation count]\label{prop:G2} \tagged{assumption + proved}
If no masked-off element is ever re-added, each element is removed at most once, so
the accepted mask moves number at most $N_u + N_x$. The actuator phase costs $N_u$
evaluations per sweep and either accepts a move or ends the phase, and likewise for
sensors, so with $C$ cycles
\begin{equation}\label{eq:evals}
  \text{evaluations} \;\leq\; N_u(A + C) + N_x(S + C) + (\text{link phase})
  \;=\; O\bigl((N_u + N_x)^2\bigr).
\end{equation}
\end{proposition}

\tagged{measured} Re-additions: $0$ on all seven runs. Sensor removals: $0$ on all
seven --- the sensor mask is never touched, and the sparsity in $N_s$ is induced by
the actuator mask and $\ell$. Cycles: exactly $2$ everywhere, which follows from
the sensor phase accepting nothing and so never perturbing the actuator fixed
point. Evaluations$/(N_u+N_x)^2$: $0.077$--$0.147$, decreasing in $n$, so
\eqref{eq:evals} is the right shape but loose --- accepted moves grow much more
slowly than $n$.

\medskip
\noindent
\tagged{not claimed} \emph{Absence of re-additions is empirical, not a consequence
of Assumption \ref{as:submod}.} Submodularity means diminishing returns, so after
removing other elements the marginal value of $u$ can only \emph{increase},
which would make re-adding it more attractive, not less.

\section{What is retired from the submitted Section IV}

\begin{center}
\begin{tabular}{ll}
\toprule
Submitted result & Reason for removal \\
\midrule
$\Phi(h)$, $h^\star$, Lemma 3 & certifies $0/150$ generations; $h_t \equiv 1 < h^\star$ \\
Definition 2(c) ($h_t$) & hop distance is a proxy for a magnitude-ordered prune \\
Lemma 4 & needs $K_{t-1}$ entrywise equal to $\Kd$ on its support \\
Propositions 1, 2 & consequences of Lemmas 3, 4; crossover is uniform, not single-point \\
Theorem 1 (drift) & single-path counting is $10^2$--$10^4$ loose; see below \\
$\mathbf{a} = \mathbf{s} = \mathbf{1}$ & assumes away the coordinate that does the work \\
\bottomrule
\end{tabular}
\end{center}

\tagged{measured} The drift route cannot be rescued by a better certificate:
substituting the \emph{exact} LQ increase for the cost bound --- the ceiling for
any certificate --- still leaves it $17.8\times$ ($5\times5$) and $196\times$
($7\times7$) below the measured decrease. The residual sits in the probability
factor $(1-p_m)^{n-c}$, which is forced: for the elite child the flipped-bit set
must contain no bit outside the certified set, because $\JEA = +\infty$ on
unstable genes and the harm of an arbitrary flip therefore has no bound.

\end{document}
